% A.74: full root and DE proof/source audits PASS; reuses the exact A71--A73 certificates.
\subsection{Odd abelian factors preserve the three two-group families}

\begin{statement}{A.74}{Proved}{Odd abelian extensions of dihedral, quaternion and semidihedral covers}
Let \(F\) be a two-power dihedral group, a generalized
quaternion two-group as in A.72, or a semidihedral two-group
as in A.73. Let \(B\) be a finite abelian group of odd order.
Both all-prime assertions of Task A.5 hold for every actual
connected finite \'{e}tale cover of \(E\setminus\{O\}\)
whose geometric Galois closure is a quotient of \(F\times B\).
Every actual intermediate, every induced cusp packet and
every permitted open derived subgroup are included.

This is a preservation theorem for an additional abelian
factor. It makes no assertion for a product of two nonabelian
Sylow factors or for arbitrary two-groups. The universal
question A.45 remains open.
\end{statement}

\begin{proof}
We apply the actual fibre-product construction of A.71,
checking its hypotheses for all three families and retaining
its separate tree and marked-path comparisons.

\paragraph{Closure groups and actual intermediate maps.}
Every subgroup of a direct product of finite groups of coprime
orders splits into its two factor intersections. For an element
\((x,y)\), Bezout powers for the two group orders produce
\((x,1)\) and \((1,y)\) in that subgroup. Applying this to
a normal quotient kernel shows that the actual closure group
has the form \(F_0\times B_0\), with \(F_0\) a quotient
of \(F\) and \(B_0\) abelian of odd order.

Every quotient \(F_0\) is again one of the three listed
types or is abelian. For a dihedral group this follows
directly by passing its rotation-and-reflection presentation
to the quotient. For a quaternion group, every nontrivial
subgroup contains its unique involution, so every proper
quotient factors through its dihedral central quotient.
For a semidihedral group its centre is exactly
\(\{1,\sigma\}\). Every nontrivial normal subgroup of a
finite \(2\)-group meets the centre nontrivially: conjugation
of the whole group on that subgroup has nonfixed orbits of
even cardinality, so its fixed subgroup has even order and
contains more than the identity. Thus every proper
semidihedral quotient also factors through its dihedral
central quotient. These are group-theoretic statements about
the actual closure, independent of a choice of source model.

Every actual covering subgroup splits as \(H_2\times H_B\).
The quotient \(B_0/H_B\) corresponds to an abelian cover of
the punctured elliptic curve. Its boundary commutator is
trivial, so it extends to a connected unramified cover of
the proper elliptic curve. After choosing an origin over
\(O\), it is an actual odd-degree elliptic isogeny
\(\eta:E'\to E\).

Let \(Z\to E\) be the actual \(H_2\)-intermediate of
the \(F_0\) closure. The function field of the given
compact curve \(Y\) is the compositum of the function
fields of \(Z\) and \(E'\) over that of \(E\). The
normal closure of the first has two-power degree, while
the second is Galois of odd degree; their intersection is
the base field. Thus the actual pullback
\[
                         Y=Z\times_E E'
\]
after smooth proper normalization is geometrically connected.
The removed points are precisely the pairs of a cusp of
\(Z\) and a kernel point of \(\eta\). We have identified
the actual cover, rather than selected an unspecified
component of a formal fibre product.

\paragraph{The precise geometric hypothesis inherited from the three families.}
If \(F_0\) is abelian, the whole given cover extends to
an elliptic isogeny and A.40 applies. More generally if
\(Z\) has genus one, so does its unramified pullback
along \(\eta\), and A.40 again applies.

For positive genus greater than one, A.70, A.72 and A.73
give the same geometric property for every actual
\(H_2\)-intermediate: after finite extension at two it
has compact-type reduction, and each selected cusp has
a distinct companion joined by rational components,
including the marked endpoint components. In those proofs
this follows from potential-good rank-zero Jacobian
reduction, actual two-primary cusp divisors and the
extended Abel-map criterion in A.41. The existence of
this rational path, rather than rank zero by itself,
is the input to the present preservation argument.

\paragraph{The actual model extends and its pullback is semistable.}
Use the actual compact-type model \(\mathcal Z/R\) just
obtained and good elliptic models \(\mathcal E,\mathcal E'\).
As in A.71, the actual map \(Z\to E\) extends to
\(\mathcal Z\to\mathcal E\): base its extended Abel map
at a cusp mapping to \(O\), and compose with the actual
extended Jacobian homomorphism. The required relative Abel
map, including the nodes, is supplied by
\cite[Theorem 1.1 and Proposition 3.2]{fghDupuyRabinoff}.
Homomorphisms between the good abelian varieties extend
over their abelian schemes. This constructs the morphism;
no general rational-map extension on a singular surface
is being assumed.

The odd-degree isogeny extends to a finite \'{e}tale
isogeny of the good elliptic models. Its finite flat
kernel has order invertible in residue characteristic two.
Therefore
\[
             \mathcal Y=\mathcal Z\times_{\mathcal E}\mathcal E'
\]
is a proper flat semistable model of the actual connected
generic curve, finite \'{e}tale over \(\mathcal Z\).
This is precisely the local model comparison in A.71.

\paragraph{Compact type is a separate consequence.}
An \'{e}tale pullback of a compact-type curve need not
have a tree dual graph. Here the actual total Galois
closure has nilpotent group \(F_0\times B_0\), and its
only possible branch point on the good elliptic base is
\(O\). The branch closure is an \'{e}tale section and
generic ramification is tame in characteristic zero.
Lehr--Matignon Theorem 6.13 therefore gives potential
good reduction of the Jacobian of this actual closure.
\cite{fghLehrMatignon2005}
The actual norm to \(Y\) makes \(J_Y\) an isogeny
quotient, hence also potentially good. After further
finite extension the above semistable model has a tree
as its dual graph. No special \(2\)-rank claim for
\(J_Y\) is needed or made.

\paragraph{The selected marked path lifts literally.}
Choose a cusp of \(Y\) above a cusp \(P\) of \(Z\).
Choose its distinct companion \(Q\) on \(Z\), and
let \(T\) be their rational joining subtree in the
geometric special fibre. Every rational component maps
constantly to the smooth elliptic curve, and the
constants agree across its nodes. The cusp sections
map to \(O\), so all of \(T\) maps to that point.
Its pullback along the finite \'{e}tale odd isogeny is
therefore literally a disjoint union of \(\deg\eta\)
copies of \(T\), indexed by the distinct special
kernel points.

The selected cusp lies on one such copy. That same copy
contains a section over \(Q\), because the odd kernel
sections have distinct reductions after making them
rational. The two generic cusps are distinct since
\(P\ne Q\). Stabilization can contract rational
components but cannot introduce a positive-genus
component into their joining path. The separately
proved tree condition makes this the unique path
required by A.42.

Thus A.42 gives the full fixed-inertia assertion at
every coefficient prime for every cusp. A.59 supplies
all normalizer assertions for the actual nilpotent
closure. Finite extensions define only the algebraic
data and models; base-point paths are used up to inner
conjugacy, and the field-restriction argument of
A.41--A.42 preserves the original packets and open
derived subgroups. This proves the statement.

No new numerical assertion is needed for this deduction.
The exact coprime-power and curve-degree support of
A.71 is retained in
\path{code/a_dihedral_abelian_products.py}; the actual
quaternion and semidihedral group checks are retained
in their respective A.72 and A.73 certificates. The
proof above verifies the quotient and model comparisons
for all group orders without extrapolating from them.
\end{proof}
